% Chapter 3 — Pipeline de procesamiento
\chapter{Pipeline de Procesamiento}

\section{Esquema general}

\begin{Verbatim}
Imágenes cámaras (Obscape API / dataset manual)
      │
      ▼
[3.1] Adquisición y preprocesado
      ├─ Normalización radiométrica
      ├─ Corrección de contraste/balance
      └─ Recorte ROI por cámara
      │
      ▼
[3.2] Segmentación semántica (SAM)
      ├─ Generación de regiones candidatas
      ├─ Mapa probabilístico arena seca / húmeda / agua
      └─ Umbral adaptativo + filtrado morfológico → máscara binaria
      │
      ▼
[3.3] Extracción de línea de costa
      ├─ Cálculo de contornos
      └─ Selección del borde húmedo-seco (polilínea en píxeles)
      │
      ▼
[3.4] Proyección homografía
      └─ (u,v) píxel → (X,Y) metros EPSG:25830
      │
      ▼
[3.5] Postprocesado + exportación
      ├─ Suavizado / simplificación métrica
      ├─ Filtros temporales (detección de anomalías)
      ├─ Cálculo de área seca (m²)
      └─ GeoJSON: ID_Camara, Timestamp, Confianza_IA, Area_Seca_m2
\end{Verbatim}

\section{Módulo offline [3.0] --- Calibración}

Ejecución única por cámara, previa al pipeline de producción.
Ver Capítulo~\ref{chap:calibration} para el detalle completo.

\begin{itemize}
  \item Transectos GNSS del IEL (GCPs en EPSG:25830)
  \item Correspondencias píxel $\leftrightarrow$ UTM
  \item Corrección distorsión de lente (parámetros intrínsecos, OpenCV)
  \item RANSAC $\to$ matriz $H$ por cámara
  \item Producto: \textbf{perfil de calibración} ($H$ + parámetros cámara + metadatos)
\end{itemize}

\section{Etapa 3.1 --- Adquisición y preprocesado}

\begin{itemize}
  \item Descarga automática via \texttt{obscape\_api.py} (cuando las cámaras
        estén registradas en la cuenta Obscape).
  \item Fase horaria prioritaria: \textbf{12:00h solar local}.
  \item Por cámara: normalización radiométrica, corrección de contraste y
        balance de blancos, recorte a la ROI definida.
\end{itemize}

\section{Etapa 3.2 --- Segmentación semántica con SAM}

Se utiliza \textbf{Segment Anything Model (SAM)} de Meta AI.

\begin{itemize}
  \item Modo \emph{zero-shot}: sin datos de entrenamiento específicos
        de Guardamar.
  \item Prompt basado en puntos semilla o bounding boxes derivados de
        las máscaras de ROI por cámara.
  \item Salida: mapa probabilístico arena seca / arena húmeda / agua.
  \item Post-SAM: umbral adaptativo + filtrado morfológico $\to$ máscara
        binaria.
\end{itemize}

\section{Etapa 3.3 --- Extracción de línea de costa}

\begin{enumerate}
  \item Cálculo de contornos sobre la máscara binaria
        (\texttt{cv2.findContours}).
  \item Selección del contorno correspondiente al borde
        \emph{arena húmeda -- arena seca}.
  \item Resultado: polilínea en coordenadas de píxel $(u, v)$.
\end{enumerate}

\section{Etapa 3.4 --- Proyección por homografía}

La homografía $H_i$ transforma coordenadas imagen a EPSG:25830:

\[
  \begin{pmatrix} X \\ Y \\ 1 \end{pmatrix}
  \sim H_i \cdot
  \begin{pmatrix} u \\ v \\ 1 \end{pmatrix}
\]

$H_i \in \mathbb{R}^{3\times3}$ se carga desde el perfil de calibración
de la cámara $i$ (ver Capítulo~\ref{chap:calibration}).

\section{Etapa 3.5 --- Postprocesado y exportación}

\begin{itemize}
  \item \textbf{Suavizado geométrico}: simplificación de la polilínea UTM
        (Douglas-Peucker o suavizado gaussiano).
  \item \textbf{Filtros temporales}: detección de anomalías respecto a la
        imagen anterior.
  \item \textbf{Cálculo de área seca}: integración del polígono en m².
  \item \textbf{Exportación GeoJSON}:
\end{itemize}

\begin{lstlisting}[language=json, caption={Estructura GeoJSON de salida (EPSG:25830)}]
{
  "type": "FeatureCollection",
  "crs": {
    "type": "name",
    "properties": { "name": "EPSG:25830" }
  },
  "features": [{
    "type": "Feature",
    "geometry": {
      "type": "LineString",
      "coordinates": [[711234.5, 4209876.3], ...]
    },
    "properties": {
      "ID_Camara":    1,
      "Timestamp":    "2026-04-30T12:00:00Z",
      "Confianza_IA": 0.93,
      "Area_Seca_m2": 14250.7
    }
  }]
}
\end{lstlisting}
